\documentclass[12pt,a4paper,landscape]{article}
\usepackage[spanish,es-nodecimaldot]{babel}
\usepackage[utf8]{inputenc}
\usepackage[T1]{fontenc}
\usepackage{lmodern}
\usepackage[top=1.1cm,bottom=1.1cm,left=1.4cm,right=1.4cm]{geometry}
\usepackage{amsmath,amssymb}
\usepackage[table]{xcolor}
\usepackage{enumitem}
\usepackage[normalem]{ulem}

\setlength{\parindent}{0pt}
\setlength{\parskip}{6pt}
\pagestyle{empty}

\newcommand{\cuerpo}{\fontsize{16.6}{20.2}\selectfont}
\newcommand{\cuerpoej}{\fontsize{19}{25}\selectfont}

\AtBeginDocument{%
  \setlength{\abovedisplayskip}{6pt}%
  \setlength{\belowdisplayskip}{6pt}%
  \setlength{\abovedisplayshortskip}{3pt}%
  \setlength{\belowdisplayshortskip}{3pt}%
  \cuerpo
}

\begin{document}

{\fontsize{20}{24}\selectfont\bfseries Prueba U} {\fontsize{17}{20}\selectfont (2 muestras indep.)}

\vspace{4pt}

\begin{minipage}[t]{0.46\textwidth}
\vspace{0pt}
Existen dos pruebas: la U y la H.

\textbf{Prueba U (Wilcoxon) - Alternativa no param. de t bimuestral}

Se quiere probar si las medias de dos muestras son significativamente iguales. Se busca probar si las muestras vienen de poblaciones idénticas (Misma distribución: misma forma, $\mu$ y varianza)

\textbf{Pasos} \hspace{1.4cm} qnorm para las absicas

1.Unimos las muestras y las ordenamos de menor a mayor. (Rango=Sum de la posic.)

2. A cada valor se le indica a que muestra pertenece (I, II) y se le asigna un rango (posición). Si dos valores son iguales, se asigna un promedio de sus rangos.

3. Sumar los rangos de cada muestra R1 y R2. se comparan y se elije el Ri menor.

4. Se define el estadístico:
\[
U_i = n_1.\, n_2 + \frac{ni.(ni+1)}{2} \; - \; Ri
\]

\begin{minipage}[t]{0.54\textwidth}
\vspace{0pt}
Si las n1 y n2 son $> 8$, Ui tiene una dist. normal. \\
gralm. prueba 2 colas
\end{minipage}\hfill
\begin{minipage}[t]{0.42\textwidth}
\vspace{0pt}
\[
z = \frac{Ui - u_{U_i}}{\sigma_{U_i}}
\]
\end{minipage}

\vspace{2pt}
\[
\mu_{U1} = \frac{n1\cdot n2}{2}
\qquad
\sigma_{U1} = \sqrt{\frac{n1\cdot n2\cdot(n1+n2+1)}{12}}
\]
\end{minipage}\hfill
\begin{minipage}[t]{0.50\textwidth}
\vspace{0pt}
\cuerpoej

\textbf{1)} \hspace{0.6cm} Se sospecha q ambos difieren en su diametro

\begin{itemize}[leftmargin=22pt,itemsep=26pt,topsep=16pt,parsep=0pt]
  \item \textbf{Arena I} ($n_1 = 15$): 0.63, 0.17, 0.35, 0.49, 0.18, 0.43, 0.12, 0.20, 0.47, 1.36, 0.51, 0.84, 0.32, 0.40, 0.45.
  \item \textbf{Arena II} ($n_2 = 14$): 1.13, 0.54, 0.96, 0.26, 0.39, 0.88, 0.92, 0.53, 1.01, 0.48, 0.89, 1.07, 1.11, 0.58.
\end{itemize}

\vspace{34pt}
\textbf{2)} \hspace{0.4cm}
\begin{minipage}[t]{0.86\textwidth}
\vspace{0pt}
\setlength{\tabcolsep}{0pt}
\renewcommand{\arraystretch}{1.7}
\begin{tabular}{@{}l@{\hspace{22pt}}l@{}}
$0.12$ (I) $\rightarrow$ Rango 1 & $0.17$ (I) $\rightarrow$ Rango 2 \\
$0.20$ (I) $\rightarrow$ Rango 4 & $0.26$ (II) $\rightarrow$ Rango 5 \\
\end{tabular}
\end{minipage}

\vspace{34pt}
\textbf{3)} \hspace{0.4cm} $\displaystyle \sum R$
\end{minipage}

\end{document}
